\SourceLeaf{23}
since the sum of them all is $=0$) that is not $\equiv 0$. Nor can the quantities
$\alpha, \beta, \gamma$ all be incongruent to one another, for then they would be congruent,
in some order, to the numbers $0, 1, 2$, so their sum would be divisible by
$3$. Among the differences $\alpha-\beta$,
$\beta-\gamma$, $\gamma-\alpha$ there is therefore exactly one that is
$\equiv 0$. Thus there are three mutually exclusive cases:

\[
\text{I.} \quad \text{If } \beta \equiv \gamma \quad(\mathrm{mod.}\ 3)
\]
\[
\text{II.} \quad \text{If } \gamma \equiv \alpha \quad(\mathrm{mod.}\ 3)
\]
\[
\text{III.} \quad \text{If } \alpha \equiv \beta \quad(\mathrm{mod.}\ 3).
\]

Which of these three cases occurs depends on $x$. Accordingly, we divide
all numbers indivisible by $p$ into three classes. We put the numbers for
which the first case occurs into the first class, those for which the second
occurs into the second, and those for which the third occurs into the third.

\textls[100]{First Theorem:} All numbers congruent modulo $p$ belong to
one and the same class.

\textls[100]{Second Theorem:} Numbers of equal magnitude and opposite signs
belong to the same class.

It is clear from this that one need only classify the numbers $1, 2, 3, \ldots,
\frac{1}{2}(p-1)$
to be able to classify all the others at once.

\textls[100]{Third Theorem:} The numbers $(3n-m)x$ and $2mx$ belong to two
consecutive classes.

\textls[100]{Fourth Theorem:} Thus if $\dfrac{2n-m}{3m} \equiv f$, so that

\[
1+f+ff \equiv 0 \quad(\mathrm{mod.}\ p), \qquad 1 \equiv f^3 \quad(\mathrm{mod.}\ p),
\]
\noindent then $x$ and the residues of $fx$ and $ffx$ belong to
three consecutive classes in reverse order.

